\section{Evaluation: Setup and Four Research Questions}
\label{sec:evaluation}
\label{sec:core-rqs}

We report an exploratory TP1 logical-demand observation and separate TP2
numerical and resident-serving records. We do not combine profiles to create
a quality, capacity, or performance result. E0 (RQ1), E3 (RQ2), the safe
blocking/async part of E1/E4 (RQ3), and integrated E5 (RQ4) retain their
historical experiment identifiers. Broad serving sweeps (E2), further
ablations (E4), and protocol overhead (E6) remain secondary evidence plans.

\subsection{Setup, controls, and evidence units}
\label{sec:eval-setup}
The TP1 middle experiment uses Qwen3.5-35B-A3B-FP8, one CUDA owner
(PP1/TP1/EP1, thread), the resident engine, lazy experts, resident KV,
context 128, batch-token limit 8, active-sequence limit 2, and prefill chunks
of 8. Its Mixed arithmetic uses BF16 activation/KV/convolution with F32
accumulation/SSM; the FP8 checkpoint is not a complete arithmetic contract.
The sealed binary is \texttt{36e6aaa5\ldots70a58c}. Table~\ref{tab:tp1-config}
reports configured limits, not physical partitions.

\begin{table}[tbp]
\centering
\small
\begin{tabularx}{\linewidth}{@{}lX@{}}
\toprule
Configured field & Value / scope \\
\midrule
Device admission / workspace & 7,168 / 2,048 MiB; not a measured total working-set cap \\
Expert cache / scratch & 64 entries; $b_{\mathrm{cache}}=335{,}544{,}320$ bytes (320 MiB, shared cache-accounting ceiling for resident, pending, and scratch); separate $b_{\mathrm{scratch}}=8$ MiB \\
KV / maximum tensor & 16 / 1,024 MiB; distinct covered limits \\
Host expert retention & Zero; not zero OS page-cache residency \\
Request queue / inflight & 2 / 2; actual workload concurrency 1 \\
Physical host limit / caches & Host cap unknown; page cache uncontrolled, not storage-cold \\
Logging / generation & \texttt{ferrule\_performance=DEBUG}; temperature 0, seed 1 \\
\bottomrule
\end{tabularx}
\caption{TP1 middle configuration. Values from the frozen case manifest;
heterogeneous or overlapping limits cannot be added as physical memory.}
\label{tab:tp1-config}
\end{table}

Here $b_{\mathrm{cache}}$ is a configured cache-accounting ceiling, not a
physical-memory bound. The plot field \texttt{cache.peak\_bytes} (raw receipt
key \texttt{expert\_owner\_peak\_bytes}) reports the lifetime high-water of
cache-charged bytes: resident plus pending-upload reservations plus scratch. That charged peak is governed by
the 320 MiB configured cache ceiling. It is not a complete physical-owner
memory peak or an independent measurement of allocator retirement; a total
physical-owner ceiling remains unknown. The configured cache flag includes
scratch according to the sealed startup plan, while $b_{\mathrm{scratch}}$
is a separate component limit. A shared cache byte ceiling constrains resident, pending,
and scratch charge; it is not a separate resident-only budget.

\paragraph{Attempt and request accounting.}
The initial CPU-prepared entry was unexecuted after its final validation was
canceled. Authorized attempt \texttt{-01} then failed before server launch
with a harness \texttt{AttributeError} for missing \texttt{subprocess.DEVNULL}:
planned 4, dispatched 0, not-run 4. No model/OOM, read, request-latency, or
quality result was obtained. Authenticated drain/post-idle checks passed,
while physical I/O and the resource contract remained unknown and GPU
quiescence unverified. This attempt is retained, not pooled with \texttt{-02}.

Attempt \texttt{-02} dispatched four serialized requests: Hello2 and the QA,
arithmetic, and Chinese canaries. All four reached successful HTTP termination;
failed and not-run counts were zero. They returned 26 usage output tokens in
30 decoder invocations. Hello2 has no reference and passed HTTP/usage only.
The other three passed their frozen visible-text/eight-output-token-usage
checks (3/3); QA is one canary, not an extended QA workload. Generated IDs,
logits, routes, and full current-model quality were not observed. This is one
independent DEBUG-instrumented run ($n=1$); observer/logging overhead was not
measured. Its request timings are exploration, not a fast baseline or
steady-state comparison.

TTFT is send-to-first-nonempty SSE content and E2E is send-to-EOF. SSE events
and usage are not exact token timestamps; no exact-token TPOT is reported.
Named request/time/log-byte anchors attribute serialized invocation deltas;
a transaction ID is not a request ID. Memory samples, logical counters,
physical traffic, and owner-ledger occupancy are distinct measurement units.
Freeze snapshot/dirty manifest, binary, model/tokenizer, workload, state
precision, backend, storage, topology, caches, and interval boundaries for
any future comparison. Offered/accepted/completed/failed/rejected/canceled/
timeout counts must remain in the denominator record.

\subsection{RQ1: Capacity and complete quality -- not established}
\label{sec:rq1}
Equations~\ref{eq:physical-capacity} and~\ref{eq:full-residency} define
deduplicated physical capacity and the full-residency test for a declared
complete weight representation;
Equation~\ref{eq:working-set} asks whether the actual streaming working set
fits each domain. The configured limits in Table~\ref{tab:tp1-config} answer
neither question alone. Table~\ref{tab:tp1-memory} reports the available memory proxies without
promoting them to continuously measured $M_d^{\mathrm{peak}}$. The
storage-demand cache figure in RQ2 is a counter/gauge reconciliation, not a
physical-working-set certification.

\begin{table}[tbp]
\centering
\small
\begin{tabularx}{\linewidth}{@{}Xrl@{}}
\toprule
Observable & Reported value & Scope \\
\midrule
Selected GPU0 sampled maximum & 4,125 MiB & 32 GPU samples \\
Server-tree PSS sampled maximum & 7,610.864 MiB & 32 complete host samples \\
Server-tree RSS sampled maximum & 7,613.770 MiB & Same sample population \\
Expert resident entries / bytes & 64 / 201,351,168 & Last invocation gauge \\
Scratch bytes & 8,388,608 & Last invocation gauge \\
Cache charged-byte high-water & 219,754,496 & Lifetime accounting gauge \\
Pending-upload / quarantine bytes & 0 / 0 & Invocation gauges only \\
\bottomrule
\end{tabularx}
\caption{TP1 \texttt{-02} sampled memory and cache-accounting gauges.
Sampled memory maxima are independent, not additive or certified physical
peaks. Cache charged-byte high-water covers resident, pending, and scratch;
its ceiling is 320 MiB, not a physical-owner bound or final ledger proof.
GPU values are selected-device aggregate telemetry, not process attribution.}
\label{tab:tp1-memory}
\end{table}

Host capacity, attributable page-cache residency, continuous peaks, and
process-owned GPU occupancy remain unknown. The OS cache was uncontrolled;
zero host expert retention does not imply zero file residency. Complete
compact weights plus scales/metadata have not been ruled out as a feasible
resident alternative. Nor do four successful HTTP responses constitute the
independent complete-output quality gate. Therefore RQ1 is open: neither
quality-qualified configured-budget feasibility nor inference beyond physical
memory/storage necessity is established by this TP1 record.

\paragraph{Later current-quality gate (2026-10-14).}
The current StorageMiddle64 full40 gate stopped before inference at
\texttt{cuInit}, with \texttt{CUDA\_ERROR\_NOT\_INITIALIZED (3)}: zero
observed numerical metrics/logits, not a numerical mismatch. Fresh-quality and
cold-owner gates were unissued; no retry occurred. Its receipt is
\path{target/validation/}\allowbreak\path{storage-middle64-current-quality-20261014-01/FINAL.json},
distinct from the earlier \texttt{-01} \texttt{DEVNULL} harness failure.
Owner initialization remains under investigation.

\paragraph{Separate TP2 numerical evidence.}
The \texttt{native-whole-capture-full40-02} comparator reports 2,178 exact
fields for each of Hello and Capital, covering prefill and first decode;
max absolute error, RMSE, and field failures are zero, with matching greedy
IDs. This is a separate full-boundary receipt from the six-row selection
below. These are bounded captured/comparator fields, including oracle-only
diagnostics that are not all physically observed device arrays, not a blanket
bitwise claim over the model. The full-boundary audit freezes this scope.
A separate multistep/lifecycle audit records the six cases in
Table~\ref{tab:tp2-numerical}: 33 calls and 35,937 exact comparator metrics.
This table deliberately selects the frozen
\path{experiments/tp2_resident/model-runtime-quality.json} receipt and its
\texttt{cuda-qwen35-mixed-bf16-tp2-local-shards-serial} profile, with test
binary \texttt{a4977b85\ldots e1cc}. A different frozen
\texttt{production-opts-01} acceptance summary, referenced by the optimized
release handoff, reports 37 calls and 40,293 exact metrics across its full40
and six-recipe selection. Its composition/model test binaries and source
snapshot differ. That summary is not substituted for the six rows here:
33/35,937 is the selected denominator for this table only, not a total of all
TP2 validation. Neither count covers the later current-quality gate above.
The frozen greedy4-02 closure records supervisor/drain/release PASS; the older
greedy4-01 supervisor failure remains a failure. These authenticated test
binaries are not public serving binaries. No serving HTTP logits were
observed, and this writing pass did not rerun payload comparisons.

\begin{table}[tbp]
\centering
\small
\begin{tabular}{@{}lrr@{}}
\toprule
TP2 test record & Calls & Comparator metrics \\
\midrule
greedy4-02 & 10 & 10,890 \\
fork2-03 & 6 & 6,534 \\
canary2-01 & 5 & 5,445 \\
recreate1-01 & 4 & 4,356 \\
cold-owner1-01 & 4 & 4,356 \\
cold-owner1-02 & 4 & 4,356 \\
\midrule
Total & 33 & 35,937 \\
\bottomrule
\end{tabular}
\caption[Separate TP2 test-binary numerical/lifecycle evidence]{Separate TP2 Mixed test-binary numerical/lifecycle evidence from
\path{experiments/tp2_resident/model-runtime-quality.json}. Exactness applies
to the frozen comparator metrics, not TP1, vLLM arithmetic, all workloads,
or the current working tree.}
\label{tab:tp2-numerical}
\end{table}

The local frozen F32 process-EP8 full-model pass also remains scoped to its
profile/snapshot, not a thread result or a capacity result. The independent
single-GPU W8A16 strict oracle gate failed on logits, ordered routes, and
internal states; argmax agreement does not reverse that failure. These
records are kept separate from TP1's three limited canaries and from the
resident timing supplement below. Neither quality nor capacity transfers
between profiles merely because a checkpoint is shared.

\paragraph{Experiment needed to close RQ1 (E0).}
Freeze representation/placement and complete-output tolerances; report each
$B_d$, $W_{\mathrm{full}}(r)$, and compact-residency alternative. Measure
physical working-set components and attribution across startup, requests,
and cleanup, stating sampling coverage. Record full-model prefill/multistep
decode for this text instance, repeated materialization/retirement, physical
request reads, retained claims, swap, and cold/warm caches. Sweep configured
budgets separately from physical capacity and retain failed/stalled points.
Use a correct bounded disk-backed reference and a host-only memory-offload
reference where feasible under the same quality/physical contract. An
infeasible reference is infeasible, not zero time or infinite speedup.

\subsection{RQ2: Request I/O -- logical demand observed, physical I/O unknown}
\label{sec:rq2}
\label{sec:tp1-storage-middle}
For request $q$, the additive invocation observables are
\begin{equation}
\begin{aligned}
 L_q&=\sum_{v\in\mathcal V_q}\Delta_v\mathrm{reader\_bytes},\\
 H_q&=\sum_{v\in\mathcal V_q}\Delta_v\mathrm{facade\_H2D\_bytes},\qquad
 \Delta_q=H_q-L_q.
\end{aligned}
\label{eq:request-io}
\end{equation}
$L_q$ is logical checkpoint-reader demand and $H_q$ is facade transfer
activity. Their difference $\Delta_q$ compares two recorded counter layers; it is neither backing-storage traffic nor a physical bus
measurement. In particular, $L_q,H_q,\Delta_q$ do not imply physical
$P_q^{\mathrm{net}}$ or $P_q^{\mathrm{blk}}$. Table~\ref{tab:tp1-storage-middle}
and Figure~\ref{fig:storage-demand-logical-transfer} apply these sums; gauge
values in Table~\ref{tab:tp1-memory} are not summed.

\begin{table}[tbp]
\centering
\small
\begin{tabularx}{\linewidth}{@{}Xrrrrr@{}}
\toprule
Request & Logical read & Uploads & Evictions & TTFT & E2E \\
 & GiB & experts & experts & s & s \\
\midrule
Hello2 & 2.813 & 960 & 896 & 4.265 & 9.186 \\
QA canary & 10.777 & 3,678 & 3,678 & 8.332 & 18.309 \\
Arithmetic canary & 10.396 & 3,548 & 3,548 & 5.308 & 14.110 \\
Chinese canary & 9.766 & 3,333 & 3,333 & 2.605 & 9.932 \\
\bottomrule
\end{tabularx}
\caption{Four serialized TP1 \texttt{-02} requests. Logical GiB
($1\,\mathrm{GiB}=2^{30}$ bytes) and HTTP timings are rounded to three
decimals; expert counts are exact invocation deltas. Not physical storage,
bus traffic, kernel timing, or comparative performance.}
\label{tab:tp1-storage-middle}
\end{table}

The 30 invocations total 69,114 checkpoint-reader calls and
36,240,064,128 logical bytes, 88,763 facade H2D copies and 36,240,248,688
facade bytes, zero recorded expert hits, 11,519 expert uploads, 11,455
evictions, and 34,557 artifact uploads. Thus the four-request aggregate is
$\Delta=184{,}560$ bytes, or 180.234375 KiB; this small difference is an
application-counter distinction, not evidence of disk traffic. Upload counts
are not observed cache misses or measured I/O credits. These counters show
repeated materialization activity, not the general registry's path integration
or physical NAS/network/block-device reads. Startup reads are outside the sums.

\paperplot{figures/storage-demand-logical-transfer.pdf}
{Paired logical checkpoint-reader and facade H2D byte deltas for four serialized
requests in storage-demand TP1 middle -02. GiB means $2^{30}$ bytes; the right
panel resolves the small H2D-minus-reader differences in KiB ($2^{10}$ bytes).
These are application counters, not physical storage I/O or physical bus
traffic; startup I/O is outside invocation accounting
(Equation~\ref{eq:request-io}).}
{fig:storage-demand-logical-transfer}

For the expert-entry count, let $I_j$ and $E_j$ be the recorded install/upload
and eviction deltas for completed invocation $j$, and let $N_j$ be the entry
gauge at that invocation boundary. The observed-boundary relation is
\begin{equation}
 N_j=N_{j-1}+I_j-E_j,\qquad j=2,\ldots,30.
 \label{eq:entry-reconciliation}
\end{equation}
For the first boundary, $N_0^{\mathrm{implied}}=N_1-I_1+E_1=0$;
this is an implied initial condition, not a measured pre-request gauge. Across
the run, 11,519 uploads minus
11,455 evictions equal the last observed 64 resident entries; all subsequent
observed boundaries have zero balance residual. This is entry-count
reconciliation, not byte conservation, continuous occupancy, or a final
owner-ledger proof.

Figures~\ref{fig:storage-demand-cache-reconciliation} and
\ref{fig:storage-demand-invocation-timeline} show the aggregate and
per-invocation forms of Equation~\ref{eq:entry-reconciliation}. The timeline
contains 3/9/9/9 invocation boundaries assigned by serialized request/log-byte
anchors, not 30 output-token or elapsed-time samples.

For the cache-gauge subset, define
\begin{equation}
\begin{aligned}
 C_j^{\mathrm{cache}}&=X_j+Q_j+S_j,\\
 C_{\mathrm{peak}}^{\mathrm{cache}}
   &=\max_{u\in\mathcal U_{\mathrm{owner}}}C_u^{\mathrm{cache}}
     \leq b_{\mathrm{cache}}.
\end{aligned}
\label{eq:cache-charge}
\end{equation}
where $X_j$ is resident expert bytes, $Q_j$ pending-upload bytes, and $S_j$
scratch bytes at invocation boundary $j$. Unknown-quarantine bytes $Z_j$ are
an already charged retained subset of $C_j^{\mathrm{cache}}$, not an
additional term. This subset is distinct from general-ledger retained $T_a$. These are cache-owner counter semantics, not evidence that
all physical allocations or grants were observed. This is an observed cache
subset, not $M_d(t)$ and not the complete $F_a+R_a+I_a+L_a+T_a$ ledger. At the
last
invocation, the recorded values are $X_j=201{,}351{,}168$ bytes,
$Q_j=0$, and $S_j=8{,}388{,}608$ bytes; the reported $Z_j$ is zero.
Equation~\ref{eq:cache-charge} adds compatible gauges at that one boundary,
not maxima or observations across time. Figure~\ref{fig:storage-demand-cache-reconciliation}
keeps the plotted gauges separate. $\mathcal U_{\mathrm{owner}}$ denotes
cache charged-state update points over the owner's lifetime, not just the 30
invocation samples. The recorded \texttt{cache.peak\_bytes} is
$C_{\mathrm{peak}}^{\mathrm{cache}}=219{,}754{,}496$ bytes, an accounting
high-water governed by \texttt{ExpertCacheLimits.max\_bytes}, namely
$b_{\mathrm{cache}}=335{,}544{,}320$ bytes. This is the
cache admission/charging contract, not independent verification of every
allocation or the general credit partition. Neither the peak nor
$b_{\mathrm{scratch}}$ certifies total physical-owner memory, continuous
$M_d(t)$, or shutdown custody closure. Unknown quarantine remains an already
charged subset and does not add another copy of the charge.

\paperplot{figures/storage-demand-cache-reconciliation.pdf}
{TP1 middle -02 entry reconciliation and cache-accounting gauges.
Uploads minus evictions ($11{,}519-11{,}455=64$) match the last observed
entries; initial zero is implied, not measured, and subsequent boundaries
reconcile. Recorded \texttt{cache.peak\_bytes}, 209.574 MiB, is the
resident-plus-pending-plus-scratch high-water, bounded by
\texttt{ExpertCacheLimits.max\_bytes} = 320 MiB, not a physical-owner bound.
Total physical-owner peak/ceiling are unobserved; no contradiction is implied.
Invocation gauges and lifetime high-waters are not accumulated or a final
ledger proof (Equations~\ref{eq:entry-reconciliation} and~\ref{eq:cache-charge}).}
{fig:storage-demand-cache-reconciliation}

\paperplot{figures/storage-demand-invocation-timeline.pdf}
{Per-invocation expert upload/eviction deltas and resident-entry observations
for the 30 completed invocations in storage-demand TP1 middle -02. Serialized
request/log-byte anchors assign the 3/9/9/9 invocations to requests. The x-axis
is log-order invocation ordinal, not elapsed time, output-token index, or
transaction-as-request ID. Entry occupancy is observed only at invocation
completion, not continuously; uploads are not inferred cache misses
(Equation~\ref{eq:entry-reconciliation}).}
{fig:storage-demand-invocation-timeline}

Startup was 18.497 s. Four separate HTTP windows sum to 51.538 s, excluding
startup and inter-request owner overhead; supervisor wall time was 75.874 s.
The origins in Figure~\ref{fig:storage-demand-timing} are explicitly
\begin{equation}
\begin{aligned}
 T_{\mathrm{start}}&=t_{\mathrm{ready}}-t_{\mathrm{launch}},\\
 T_q^{\mathrm{first}}&=t_q^{\mathrm{first\ content}}-t_q^{\mathrm{send}},\qquad
 T_q^{\mathrm{end}}=t_q^{\mathrm{EOF}}-t_q^{\mathrm{send}}.
\end{aligned}
\label{eq:timing-origins}
\end{equation}
These are distinct denominators. Startup has a launch origin; each request
bar has its own HTTP send origin, not one cumulative timeline. It therefore
does not imply an end-to-end throughput
rate or an I/O bottleneck share. $P_q^{\mathrm{net}}$ and
$P_q^{\mathrm{blk}}$ remain unknown. Logical reader/syscall bytes can be
served by cache; facade transfers are not physical bus measurements.

\paperplot{figures/storage-demand-timing.pdf}
{Launch-to-ready startup and per-request HTTP TTFT/E2E in the single successful
DEBUG-instrumented storage-demand TP1 middle -02 run ($n=1$). TTFT is send to
first nonempty SSE content; E2E is send to EOF. Startup has a different origin
and is not part of the request bars. Instrumentation overhead is unmeasured;
there are no statistics, confidence intervals, kernel timings, or comparative
speedups (Equation~\ref{eq:timing-origins}).}
{fig:storage-demand-timing}

To close RQ2 (E3), align application, process, mount/network, and backing-block
observations to the same requests and declared cache policy, with foreign
traffic and attribution limits. Distinguish startup from requests, runtime
cache from OS/shared file cache and storage-server cache, and physical disk
reads from memory migration. A warm request with no backing read is valid but
not evidence of request-time disk access. Swap traffic must be disclosed,
not silently counted as the intended storage path.

\subsection{RQ3: Safe scheduling -- reference not implemented}
\label{sec:rq3}
The same-contract safe-blocking reference on the actual TP1 path is
\textbf{not implemented}. Neither \texttt{-02}'s HTTP latencies nor the
resident TP2/vLLM supplement compare blocking and asynchronous scheduling.
There is no measured physical-overlap, event-driven-progress, or causal
scheduling benefit in the current evidence.

The E1/E4 experiment must first audit the actual path specified in
Section~\ref{sec:implementation}, including bypasses. Implement a correct
bounded blocking reference and serialize selected stages only while keeping
weights, arithmetic, quality, caches, ledgers, sharing/custody, and publication
rules fixed. Unsafe buffer reuse or unbounded queues are not simplification
controls. Record submission/completion and true provider/accelerator intervals,
operation/generation identities, dependency readiness, queue/claim occupancy,
blocked workers, locks, and synchronization. Correlate clocks and report
uncertainty before intersecting storage/transfer/compute intervals. Host
enqueue intervals alone cannot prove physical overlap. Include exhausted
budgets, stalled owners, and full queues to expose cyclic waits/starvation
under the stated conditional-liveness assumptions.

\subsection{RQ4: Custody and failure -- final resource contract unknown}
\label{sec:rq4}
For TP1 \texttt{-02}, authenticated descendant drain reached \texttt{ECHILD};
wrapper post-idle and independent all-eight-GPU observational idle checks
passed. Nevertheless \texttt{resource\_contract\_complete=unknown} and the
supervisor raw \texttt{gpu\_quiescence=unverified} remain unchanged. Invocation
pending-upload/quarantine gauges of zero are not a final owner ledger.
Process drain and idle telemetry establish neither device fences nor source
completion nor the partition in Equation~\ref{eq:credit-conservation}.

Algorithm~\ref{alg:terminal-custody} therefore specifies the required
behavior, not the measured closure of this run. E5 must inject delayed/stale
completion, stage-wise cancellation, primary-plus-cleanup failure, owner loss,
and lost ACKs on the integrated path. Track exact publication/reuse, original
versus repeated submissions, ownership transfers, retained $T_a$, named
completion proofs, drain, and post-failure service. Reconcile grants and
allocations before and after proven quiescence; distinguish deterministic
protocol injection from real hardware failure. Unknown should remain charged
and unusable without fabricated rollback or recovery. A failed/stalled
provider also tests liveness assumptions, not just safety. Integrated fault
conservation, eventual reclamation, and final TP1 custody closure remain open.

\subsection{Supplement: resident observations, not an RQ1--RQ4 substitute}
\label{sec:resident-supplement}
\paragraph{Supplementary resident TP2 observations.}
The frozen route-batch crossblocks record contains two independent blocks
in F1--V1 and V2--F2 order, with four fresh server owners. Each slot excludes
three QA8 requests and then measures five requests, each with five input
and eight output tokens: 40 output tokens per measurement window.
Table~\ref{tab:resident-crossblocks} transcribes these existing observations;
no new measurement is introduced here.

\begin{table}[t]
\centering
\small
\begin{tabular}{@{}lrrrrrr@{}}
\toprule
Slot & Output & TTFT & E2E & Startup & GPU & PSS \\
 & tok/s & ms & ms & s & GiB/card & GiB \\
\midrule
F1 & 14.317 & 176.248 & 523.399 & 349.585 & 21.04 & 61.05 \\
V1 & 10.376 & 132.676 & 694.228 & 66.239 & 19.28 & 7.71 \\
V2 & 10.330 & 134.651 & 699.213 & 62.814 & 19.28 & 7.32 \\
F2 & 14.804 & 171.333 & 515.106 & 366.133 & 21.04 & 60.14 \\
\bottomrule
\end{tabular}
\caption{Descriptive resident TP2 crossblocks ($n=2$). TTFT and E2E are
request means; GPU and server-tree PSS are sampled peaks, not exact peaks
or per-owner GPU attribution. F denotes Ferrule; V denotes vLLM.}
\label{tab:resident-crossblocks}
\end{table}

Ferrule has a higher output rate and lower mean E2E in both short windows;
vLLM has lower mean TTFT, shorter startup, and lower sampled GPU and host
PSS. The descriptive F/V rate ratios are 1.380 and 1.433, not statistical or
causal speedups. The record uses a pinned checkpoint and tokenizer with
TP2, context 128, batch 8, sequence limit 2, and concurrency 1. Ferrule's
FP8-checkpoint/BF16 activation, KV and convolution/F32 SSM path uses
HF-eager-qualified GemmEx/routing, whereas vLLM uses BF16/FP8 Marlin.
These are not bit-exact or arithmetic-equivalent implementations, and the
record retains \texttt{fairmatch=false} and \texttt{exact\_speedup=null}.
Ferrule's 22 GiB/device admission budget, including a 4 GiB workspace with
a 3 GiB arena, is not identical to vLLM's utilization setting of 0.8.
The common sampled 22 GiB ceiling is not a hard physical partition or a
matched total GPU/host working-set budget.

The measured vLLM deployment is \texttt{0.28.0+cu129}, with
\texttt{Torch 2.13.0+cu129}; its installed identities are frozen under
\path{target/vendor/vllm028-py312}, not inferred from a 0.18 source checkout.
Visible-text/usage canaries and Ferrule's separate scoped oracle records do
not establish cross-engine internal-state parity or universal task quality.
These resident windows neither demonstrate request-time physical disk reads
nor establish host-memory insufficiency, storage necessity, or E0 acceptance.
The immutable source record is
\path{papers/arxiv/experiments/tp2_resident/runs/route-batch-vllm-crossblocks-20261012-01/crossblocks-results.md},
with raw metrics in the adjacent \path{crossblocks-results.json} and frozen
identities in \path{sealed-identities.json}. Historical pilots and failed
attempts remain separate records, not pooled repetitions of this configuration.


\paragraph{Supplementary startup-metadata observations.}
The sealed release \texttt{36e6aaa5\ldots70a58c} combines prepared admission
and shared \texttt{Arc<Path>} checkpoint read-plan metadata. Its strict TP2
S1/S2 smoke lifecycles and ordinary F1 completed with source-bound executable
and mapped-library checks. S1/S2 passed visible greedy canary, reuse,
cancellation-policy and post-cancel gates, with authenticated descendant
drain. These are scoped serving checks, not a fresh 14-million-logit numerical
validation or proof of device quiescence from process exit.
Table~\ref{tab:startup-metadata} separates smoke and performance records.

\begin{table}[t]
\centering
\small
\begin{tabularx}{\linewidth}{@{}Xrrr@{}}
\toprule
Record & Startup & Peak tree PSS & Usage rate \\
 & s & GiB & tok/s \\
\midrule
36e6 S1 (smoke) & 222.079 & 42.716 & --- \\
36e6 S2 (smoke) & 139.855 & 44.416 & --- \\
36e6 F1 (ordinary) & 140.052 & 43.573 & 13.899 \\
e674 F1 (historical ordinary) & 351.114 & 60.149 & 14.227 \\
\bottomrule
\end{tabularx}
\caption{Startup and sampled PSS observations. Smoke slots have no QA8
throughput measurement. The historical ordinary row is context, not a
paired before/after control.}
\label{tab:startup-metadata}
\end{table}

F1 excludes three QA8 requests and measures five, each with five input and
eight output tokens: 40 usage output tokens over 2.877869 s. Its request
TTFT mean/p50/p95 are 192.944/187.513/209.700 ms, descriptive summaries of
five requests, not independent-trial confidence intervals. S1, S2, and F1
are not pooled, and the best startup is not treated as a stable constant.
The historical ordinary e674 F1 is a different series from the crossblocks
in Table~\ref{tab:resident-crossblocks}; no fresh BEFORE or de288 control
was run. Page-cache state is uncontrolled, not cold. Startup, sampled PSS,
and output-rate tradeoffs therefore do not identify a causal optimization
effect, live-heap reduction, freed OS memory, or a storage-capacity result.

Five F1 proc snapshots, in first-identification, ready, post-warmup,
post-measurement, and post-exit order, cost
0.643/623.616/725.908/625.392/0.085 ms to read and validate. Their recorded
bounds all lie outside the measured window; receipt writes are also outside,
but persistence cost was not separately measured. This excludes those reads
from the stated denominator, not observer overhead from the system. On the
recorded Linux 4.19 deployment, \texttt{Pss\_Anon}/\texttt{Pss\_File} are
unavailable; the post-exit snapshot is unavailable, not zero memory.

The source-relation proof records unchanged model mathematics, backend/native
kernels, and runtime, while checkpoint representation differs. CPU IO-byte,
ordering, and source-check tests support compatibility, not fresh GPU numerical
acceptance; full-checkpoint payload hashing was not performed. Thus metadata
representation and source-bound plans have implementation evidence, while the
physical request-time storage frontier remains unvalidated.

The local reproducibility entry \path{REPORT.md} is in:
\begin{quote}
\small\path{target/validation/startup-metadata-owner-20261013-01}
\end{quote}
Its \path{RESULTS.json} supplies the full binary hash and raw-pin manifest;
\path{executed-source-pins.json} identifies executed Python/config revisions.
The F1 raw HTTP and \path{proc-timeline.json} are under the repository's
paper experiment runs directory, at:
\begin{quote}
\small
\path{papers/arxiv/experiments/tp2_resident/runs}/\\
\path{startup-metadata-info-after-20261013-01}/\\
\path{runs/F1/session}
\end{quote}
Historical ordinary raw pins are in \path{routebatch-performance.json},
under the sibling directory \path{route-batch-info-after-20261012-01}.
These working-directory artifacts still require portable archival; they have
not been moved or repackaged for publication.


\subsection{Evidence gates and reproducibility}
\begin{table}[tbp]
\centering
\small
\begin{tabularx}{\linewidth}{@{}l>{\raggedright\arraybackslash}X>{\raggedright\arraybackslash}X@{}}
\toprule
Question & Available observation & Gate still open \\
\midrule
RQ1 & TP1 HTTP/canaries and sampled device/tree memory & Complete numerical quality, physical capacities/working-set peaks, compact-residency alternative \\
RQ2 & 36,240,064,128 logical bytes; $L/H$ transfer difference and entry/cache plots & Request physical network/block I/O and controlled caches \\
RQ3 & Algorithms and path-audit requirements & Safe-blocking reference not implemented; no physical overlap trace \\
RQ4 & Authenticated drain and observational idle & Resource contract unknown; device quiescence unverified; integrated fault/ledger evidence \\
TP2 supplement & Separate numerical tests; two resident crossblocks & No TP1 transfer, arithmetic parity, fair speedup, or capacity proof \\
\bottomrule
\end{tabularx}
\caption{Observable-to-claim map. An open gate is not filled by a result from
another representation, binary, topology, or workload.}
\label{tab:evidence}
\end{table}

The fixed TP1 reports are relative to \path{papers/arxiv}, under
\path{experiments/storage_demand/runs/}:
\begin{quote}
\small
\path{storage-demand-tp1-middle-20261013-01/}\allowbreak\path{report.json}\\
\path{storage-demand-tp1-middle-20261013-02/}\allowbreak\path{report.json}
\end{quote}
The latter has adjacent \path{case.json}, \path{server.log},
\path{samples.jsonl}, \path{supervisor.json}, \path{custody.json}, and
\path{post-idle.json}. The independent closeout is under
\path{../../target/validation/}:
\begin{quote}
\small
\path{storage-demand-middle-gpu-20261013-02/}\allowbreak\path{closure.json}
\end{quote}
The deterministic offline receipt analysis is in
\path{experiments/storage_demand/analysis/output/}; it reports source
consistency PASS with 41 bound inputs, 39 consumed receipt pins, and zero
divergences. That is a receipt-consistency check, not a new GPU or numerical
acceptance result. The TP2 numerical sources are
\path{experiments/tp2_resident/full-boundary-audit.json} and
\path{experiments/tp2_resident/model-runtime-quality.json}. The separate
37-call selection and its handoff are under \path{../../target/validation/}:
\begin{quote}
\small
\path{tp2-mixed-production-opts-01/}\allowbreak\path{acceptance-summary.json}\\
\path{tp2-mixed-optimized-release-20261009/}\allowbreak\path{handoff.json}
\end{quote}
All paths in this paragraph are relative to \path{papers/arxiv}.

Publication requires portable frozen artifacts: source plus dirty-state
manifest, binary, model, and input-processing identities, commands and environment,
raw results, traces/clock-correlation method, retained failures, and
deterministic plot/table generation. Local ignored \path{target/} files and
private transcripts are not public artifacts. Do not silently omit failed
trials. Before confirmatory collection, freeze trial/block units, repetition
counts, order, stopping rules, contrasts, and interval methods. Windows in one
server are not independent repetitions. The historical minimum of five
independent paired performance blocks applies if performance is promoted to
a primary claim; it neither specifies the core RQ repetition counts nor
guarantees power. No statistical victory or universal quality threshold is
inferred from these exploratory records.
